\documentclass[conference]{IEEEtran}
\IEEEoverridecommandlockouts

\usepackage{cite}
\usepackage{amsmath,amssymb,amsfonts}
\usepackage{algorithmic}
\usepackage{graphicx}
\usepackage{textcomp}
\usepackage{xcolor}
\usepackage{booktabs}
\usepackage{multirow}
\usepackage{listings}
\usepackage{adjustbox}
\usepackage{threeparttable}
\usepackage{url}
\usepackage{hyperref}

\lstset{
    basicstyle=\ttfamily\scriptsize,
    keywordstyle=\color{blue},
    commentstyle=\color{gray},
    stringstyle=\color{red},
    numbers=left,
    numberstyle=\tiny,
    breaklines=true,
    showstringspaces=false
}

\def\BibTeX{{\rm B\kern-.05em{\sc i\kern-.025em b}\kern-.08em
    T\kern-.1667em\lower.7ex\hbox{E}\kern-.125emX}}

\hyphenation{
  Glo-bal-Ave-rage-Pool-ing
  Glo-bal-Avg-Pool-ing
  Swift-UI
  ca-te-go-ri-cal
  Wire-less
  Neu-ral
  ar-chi-tec-ture
  ar-chi-tec-tures
  ar-chi-tec-tur-al
  clas-si-fi-ca-tion
  clas-si-fi-er
  in-fer-ence
  pre-pro-cess-ing
  pre-pro-cessed
  de-ploy-ment
  de-ploy-able
  pa-ra-me-ters
  pa-ra-me-ter
  per-for-mance
  com-plex-i-ty
  evo-lu-tion-ary
  op-ti-mi-za-tion
  con-vo-lu-tion-al
  con-vo-lu-tions
  re-cog-ni-tion
  com-mu-ni-ca-tion
  com-pu-ta-tion-al
  con-nec-tiv-i-ty
  pro-pa-ga-tion
  spec-trum
  spec-tral
  mo-bile
  ac-knowl-edged
  lim-i-ta-tions
  mi-cro-elec-tron-ics
  ac-cel-er-a-tor
  ac-cel-er-a-tors
  quan-ti-za-tion
  charac-ter-i-za-tion
  char-ac-ter-is-tics
  ad-dress-ing
  elec-tro-mag-net-ic
  meth-od-ol-o-gies
  straight-for-ward
  nor-mal-i-za-tion
  reg-u-lar-i-za-tion
  seg-men-ta-tion
  si-mul-ta-ne-ous-ly
  down-con-ver-sion
  im-per-fec-tions
  im-prac-ti-cal
  pre-dict-able
  gen-er-al-i-za-tion
  trans-mis-sions
  challen-ges
}

\begin{document}

\title{Hardware-Efficient Neural Architecture Search for Wireless Signal Classification on Edge Devices}

\author{
\IEEEauthorblockN{Félix D. Suárez Bonilla, Jose M. de la Rosa, G. Liñán-Cembrano}
\IEEEauthorblockA{Instituto de Microelectrónica de Sevilla, IMSE-CNM (CSIC/Universidad de Sevilla), Spain \\
Email: felix.suarez@imse-cnm.csic.es}
\thanks{The authors acknowledge the limited use of artificial intelligence tools to assist with manuscript preparation, specifically for writing assistance and style revision. All content has been thoroughly reviewed, validated, and remains the full responsibility of the authors.}
\thanks{This work was supported by the European Union's Recovery and Resilience Facility-Next Generation, in the framework of the General Invitation of the Spanish Government's public business entity Red.es; by PID2022-138078OB-I00 and PDC2023-145808-I00, funded by MICIU/AEI/10.13039/501100011033; and by grant USECHIP (TSI-069100-2023-001), funded by the Secretary of State for Telecommunications and Digital Infrastructure, Ministry for Digital Transformation and Civil Service, and by the European Union–NextGenerationEU, by ERDF A way of making Europe.}
}

\maketitle

\begin{abstract}
This paper presents an evolutionary Neural Architecture Search (NAS) methodology for hardware-efficient wireless signal classification on edge devices. The updated NAS configuration discovered an ultra-compact model with 3,539 parameters that classifies LTE, DVB-T, and WiFi signals with 90.1\% test accuracy on real-world over-the-air data. The inference was implemented in an iOS application and tested on actual iPhone hardware.
\end{abstract}

\begin{IEEEkeywords}
Neural Architecture Search, Wireless Signal Classification, Edge Computing, Hardware-Efficient Deep Learning, Core ML, Embedded Systems, LTE, DVB-T, WiFi
\end{IEEEkeywords}

% =========================================================
\section{Introduction}
% =========================================================

Automatic wireless signal classification has become essential for spectrum monitoring, cognitive radio networks, and interference detection \cite{oshea2017introduction}. The growing complexity of electromagnetic environments demands classification methodologies deployable on resource-constrained edge devices such as mobile SoCs, FPGAs, and ASICs \cite{gurbuz2022deep}. Traditional approaches relying on manual feature engineering face computational constraints on such platforms, while complex deep learning models exceed practical memory and power budgets.

Neural Architecture Search (NAS) addresses these challenges by automatically discovering architectures tailored for target hardware \cite{zoph2016neural, real2019regularized}. This work presents the complete design flow from algorithm development to hardware deployment of an efficient wireless signal classifier on edge devices, with a practical deployment example using Apple's Core~ML framework \cite{coreml_guide}. Without loss of generality, we classify LTE, DVB-T, and WiFi using real-world over-the-air IQ signals from Fontaine and Shahid's dataset \cite{fontaine2023dataset, fontaine2019low, fontaine2023github}.

The architecture \ref{fig:architecture} ,consisting of only two convolutional layers followed by global average pooling and a small dense head, facilitates straightforward hardware mapping with predictable, streaming memory access patterns on FPGAs and ASICs, since each layer processes data sequentially without skip connections or complex branching \cite{gurbuz2022deep}.

Our contributions include an evolutionary NAS algorithm jointly optimizing accuracy and compactness, an ultra-compact model with only 3,539 parameters achieving 90.1\% test accuracy, Core ML deployment with sub-millisecond inference on iPhone Neural Engine, a complete application for iOS, and experimental validation on mobile hardware \cite{fontaine2019low, liu2017wireless}.

To support reproducibility, the complete implementation, training scripts, and updated artifacts are publicly available at \url{https://shorturl.at/0rjX3}.

% =========================================================
\section{Related Work}
% =========================================================

\subsection{Signal Classification and Hardware Complexity}

Table~\ref{tab:comparison_modern_wireless} summarizes state-of-the-art wireless signal classification methods with emphasis on model complexity. Zhang et al. \cite{zhang2021signal} achieved 98.5\% on synthetic and 96.2\% on over-the-air data using CNN-RNN architectures. Ide et al. \cite{ide2024evaluation} obtained 96.8\% with ResNet \cite{he2016deep} ($\sim$138k parameters). Shebert et al. \cite{shebert2021wireless} attained 94.5\% but with 2.1M--67.1M parameters, impractical for edge deployment. Fontaine et al. \cite{fontaine2019low} achieved 97.8\% with 213k parameters and 88.0\% with 6.4k parameters on the same dataset we employ. While MobileNets \cite{howard2017mobilenets} and compression techniques \cite{han2015deep} have shown promise for mobile vision, existing signal classification methods suffer from either excessive complexity or lack of edge deployment validation.

\begin{table*}[!t]
\centering
\caption{Comparison of Signal Classification Methods for Modern Wireless Technologies with Hardware Complexity Analysis$^{a,b,e}$}
\resizebox{\textwidth}{!}{%
\begin{threeparttable}
\begin{tabular}{|l|l|c|c|c|c|}
\hline
\textbf{Model} & \textbf{Accuracy} & \textbf{\# Classes} & \textbf{Technologies} & \textbf{Dataset$^{a}$} & \textbf{Model Complexity$^{b}$} \\
\hline
\textbf{This Work (NAS)} & \textbf{90.1\%} & \textbf{3} & \textbf{LTE, WiFi, DVB-T} & \textbf{OTA} & \textbf{Very Low (3.5k params)} \\
\hline
Ide et al. 2024 \cite{ide2024evaluation} & 96.8\% & 3 & 5G, LTE, WLAN & SYN & Low ($\sim$138k params) \\
\hline
Shebert et al. 2021 \cite{shebert2021wireless} & 94.5\% at SNR $>$ 0 dB & 6 & LTE, 5G NR, WiFi 6, BLE & SYN & High (2.1M -- 67.1M) \\
\hline
Zhang et al. 2021 \cite{zhang2021signal}$^{e}$ & 98.5\% (SYN), 96.2\% (OTA) & 3 & WiFi, LTE LAA, 5G NR-U & OTA \& SYN & Low ($\sim$165k params)$^{e}$ \\
\hline
Fontaine et al. 2019 \cite{fontaine2019low} & 97.8\% / 88.0\% & 3 & LTE, WiFi, DVB-T & OTA & Medium (213k) / Very Low (6.4k) \\
\hline
Gu et al. 2019 \cite{gu2019deep} & 95\% at high SNR & 2 & LTE-U, WiFi & OTA & Low (99.9k params) \\
\hline
\end{tabular}
\begin{tablenotes}
\item[a] OTA: Over-the-Air. SYN: Synthetic.
\item[b] Complexity categorization: Very Low ($<$20k), Low (20k--500k), Medium (500k--5M), High ($>$5M params).
\item[e] Estimated from the reported architecture (CNN-LSTM) due to missing hyperparameter details in the original work.
\end{tablenotes}
\end{threeparttable}
}
\label{tab:comparison_modern_wireless}
\end{table*}

\subsection{Neural Architecture Search}

NAS automates the discovery of optimal neural network architectures \cite{elsken2019neural}. Zoph and Le \cite{zoph2016neural} introduced reinforcement learning-based NAS, while Real et al. \cite{real2019regularized} showed that evolutionary algorithms achieve comparable results at lower cost.

We adopt an evolutionary approach because it naturally handles discrete architectural choices---such as the number of layers, filter counts, and kernel sizes---and allows flexible integration of hardware-aware constraints directly into the fitness function \cite{real2019regularized}. Unlike gradient-based NAS methods that require continuous relaxations of the search space, evolutionary strategies evaluate complete architectures through actual training runs, producing directly deployable models without a separate architecture derivation step.

Our implementation uses tournament selection, uniform crossover of architectural genes, and a 10\% mutation rate to explore the search space efficiently. We use a scalarized fitness function combining classification error and a parameter-count penalty into a single minimization objective (detailed in Section~III-C), simplifying the search while enforcing the desired compactness--accuracy trade-off  \cite{Tournament2026}.

% =========================================================
\section{System Architecture and Methodology}
% =========================================================

%\subsection{Signal Processing Pipeline}

The RF signal reception chain (Fig.~\ref{fig:rf_chain}) captures electromagnetic signals through an antenna, filters them via an RF bandpass filter, and amplifies them with a Low-Noise Amplifier (LNA). The signal then enters a quadrature downconversion stage producing the In-phase (I) and Quadrature (Q) components. Each branch passes through Low-Pass Filters (LPF) before dual analog-to-digital converters (ADCs) simultaneously digitize the I and Q signals. Finally, we reach the AI engine that hosts the CNN model.

\begin{figure*}[!t]
\centering
\includegraphics[width=1.0\textwidth]{Images/rf_chain.png}
\caption{RF reception signal chain: antenna, filtering, LNA, quadrature I/Q demodulation with PLL synthesizer, baseband filtering, dual ADC, and DSP with AI-based classification.}
\label{fig:rf_chain}
\end{figure*}

\subsection{Dataset and Preprocessing}

We formulate wireless signal classification as predicting $y \in \{\text{LTE}, \text{DVB-T}, \text{WiFi}\}$ from input IQ signal $\mathbf{x} \in \mathbb{R}^{512 \times 2}$, subject to constraints of fewer than 20,000 parameters ($<$80~KB in FP32), inference latency below 1~ms, and regular structure amenable to hardware accelerator mapping.

We employ the Fontaine and Shahid dataset \cite{fontaine2023dataset}, containing over-the-air IQ signals from LTE, WiFi, and DVB-T transmissions collected at six locations in Gent, Belgium, using USRP equipment at 1~Msps \cite{fontaine2023github}. The dataset provides authentic signal characteristics with naturally varying strength and noise levels \cite{fontaine2019low}, ensuring robustness to real imperfections. The raw corpus comprises 135 binary files (45 per class), which are then segmented into 512-sample IQ chunks used as training/validation/test samples in our NAS pipeline.

The preprocessing pipeline applies per-channel statistical normalization, requiring only a single multiply-accumulate operation per sample with precomputed $1/\sigma$ coefficients. Signals are segmented into non-overlapping 512-sample chunks, where the power-of-two length facilitates efficient memory addressing in hardware. For NAS experiments, we use balanced subsets per class for train/validation/test and reshape segments to $(N, 512, 2)$ for 1D convolutional processing.

\subsection{Hardware-Aware Evolutionary NAS}

Our updated evolutionary NAS uses a population of 16 architectures evolved over 10 generations, with 8 training epochs per architecture during search. The search space focuses on compact 1D convolutional models (including standard and separable convolutions), optional small LSTM blocks, compact dense layers, and pooling options including global average pooling. The fitness function is tuned to prioritize accuracy while softly penalizing large models:

\begin{equation}
F = 1.2 \cdot (1-A) + 0.25 \cdot P_{\mathrm{penalty}}
\label{eq:fitness}
\end{equation}

\noindent where $A \in [0,1]$ is validation accuracy and $P_{\mathrm{penalty}}$ is a piecewise parameter penalty with mild penalty up to $\sim$10k parameters and stronger penalty beyond $\sim$16k parameters. Each generation proceeds through evaluation, elite selection, tournament selection, crossover, mutation, and elitist replacement.

\begin{figure*}[!t]
\centering
\includegraphics[width=0.8\textwidth]{Images/architecture_diagram_SMACD.png}
\caption{NAS-generated compact architecture for IQ classification: Conv1D(16)-Conv1D(32)-GlobalAveragePooling-Dense(16)-Dense(3). Total: 3,539 parameters, $\sim$13.8~KB (FP32) or $\sim$3.5~KB (INT8).}
\label{fig:architecture}
\end{figure*}

\subsection{Discovered Architecture}

The updated NAS run discovered a compact architecture (Table~\ref{tab:architecture} and Fig.~\ref{fig:architecture}) consisting of two 1D convolutional layers (both with kernel size 5), batch normalization, dropout, global average pooling, and a compact dense head (16 and 3 units). Training uses Adam \cite{kingma2014adam} with categorical cross-entropy and dropout regularization \cite{srivastava2014dropout}. The final model contains 3,539 parameters (3,443 trainable and 96 non-trainable).

\begin{table}[htbp]
\caption{NAS-Generated Architecture with Layer-Level Details}
\begin{center}
\resizebox{\columnwidth}{!}{%
\begin{tabular}{lccccc}
\toprule
\textbf{Layer Type} & \textbf{Output Shape} & \textbf{Params} & \textbf{MACs} & \textbf{Memory$^{a}$} \\
\midrule
Input & (512, 2) & 0 & 0 & 4.0 KB \\
Conv1D-1 (k=5, 16 filters) & (512, 16) & 176 & 81,920 & 0.69 KB \\
BatchNorm-1 + Dropout & (512, 16) & 64 & -- & 0.25 KB \\
Conv1D-2 (k=5, 32 filters) & (512, 32) & 2,592 & 1,310,720 & 10.13 KB \\
BatchNorm-2 + Dropout & (512, 32) & 128 & -- & 0.50 KB \\
GlobalAveragePooling1D & (32) & 0 & 16,384$^{b}$ & -- \\
Dense-1 (16 units) + Dropout & (16) & 528 & 512 & 2.06 KB \\
Dense-2 (3 units, softmax) & (3) & 51 & 48 & 0.20 KB \\
\midrule
\textbf{Total} & & \textbf{3,539} & \textbf{1,409,584} & \textbf{$\sim$13.8 KB} \\
\bottomrule
\end{tabular}%
}
\label{tab:architecture}
\end{center}
\vspace{-2mm}
\footnotesize{$^{a}$Weight memory in FP32 (4 bytes/param). INT8 quantization reduces to $\sim$3.5~KB. MACs: Multiply-Accumulate operations per inference. $^{b}$Approximate reduction operations for temporal averaging.}
\end{table}

The compact convolutional pipeline produces predictable streaming memory access patterns suitable for hardware accelerators. The model fit easily within typical FPGA BRAM resources (e.g., a Xilinx Artix-7 XC7A35T provides 225~KB), leaving significant memory headroom for buffering and control.

%\subsection{iOS Deployment via Core~ML}

For the deployment, the trained TensorFlow model \cite{abadi2016tensorflow} is converted to Core~ML format \cite{coreml_guide} using \mbox{coremltools}, targeting the \mbox{mlprogram} format for Neural Engine optimization. The conversion removes dropout layers, fuses batch normalization into convolutional weights, and adapts the data layout for efficient mobile inference. The compact model size ensures the network can reside in local memory, and the iOS application executes the complete pipeline; IQ loading, normalization, segmentation, and inference, entirely on-device without cloud connectivity.

% =========================================================
\section{Experimental Results}
% =========================================================

The evolutionary algorithm converged over 10 generations, with best fitness improving from 0.1912 to 0.1344 (lower is better). During search, validation accuracy reached 88.8\% with 3,539 parameters, showing stable convergence toward compact, hardware-friendly architectures.

The model achieved 94.5\% training, 87.7\% validation, and 90.1\% test accuracy, showing a strong compactness-accuracy compromise for strict edge constraints. Table~\ref{tab:class_metrics} reports per-class performance with macro-averaged precision of 90.2\%, recall of 90.1\%, and F1-score of 90.1\%. The dominant confusion remains between DVB-T and WiFi, likely due to spectral similarities between OFDM-based technologies at baseband.

\begin{table}[htbp]
\caption{Per-Class Classification Metrics and Confusion Matrix (\%)}
\begin{center}
\resizebox{\columnwidth}{!}{%
\begin{tabular}{lcccccc}
\toprule
 & \multicolumn{3}{c}{\textbf{Classification Metrics}} & \multicolumn{3}{c}{\textbf{Confusion Matrix}} \\
\cmidrule(lr){2-4} \cmidrule(lr){5-7}
\textbf{Class} & \textbf{Prec.} & \textbf{Rec.} & \textbf{F1} & \textbf{LTE} & \textbf{DVB-T} & \textbf{WiFi} \\
\midrule
LTE & 89.2\% & 92.8\% & 91.0\% & 92.8 & 6.4 & 0.8 \\
DVB-T & 85.2\% & 86.6\% & 85.9\% & 10.6 & 86.6 & 2.8 \\
WiFi & 96.2\% & 90.8\% & 93.4\% & 0.6 & 8.6 & 90.8 \\
\midrule
\textbf{Macro Avg} & \textbf{90.2\%} & \textbf{90.1\%} & \textbf{90.1\%} & & & \\
\bottomrule
\end{tabular}%
}
\label{tab:class_metrics}
\end{center}
\end{table}

The TensorFlow-to-Core~ML deployment path is maintained for edge execution. Fig.~\ref{fig:app} shows the iOS application interface used for real-time classification. Experimental deployment tests were conducted on an iPhone 17 Pro Max featuring an Apple A19 Pro SoC (3~nm, 6-core CPU with performance and efficiency cores, and improved GPU), 12~GB LPDDR5X RAM, and 256~GB storage. For the updated 3,539-parameter model, the current work focuses on algorithmic optimization and offline evaluation; detailed on-device latency/power characterization is left as immediate future work.

\begin{figure}[htbp]
\centering
\includegraphics[width=0.55\columnwidth]{Images/app_interface_SMACD.png}
\caption{iOS application showing real-time wireless signal classification with the NAS-optimized model on the Neural Engine.}
\label{fig:app}
\end{figure}

%All baselines were trained on the same Fontaine dataset \cite{fontaine2023dataset} with identical preprocessing and training procedures. Table~\ref{tab:comparison} summarizes the results. Our updated model achieves 90.1\% accuracy with 3,539 parameters---a 98.6\% reduction relative to a 250k-parameter traditional CNN baseline, while still enabling fully on-chip SRAM operation.

%Compared with Fontaine et al. \cite{fontaine2019low}, their IQ-CNN achieves 97.8\% but requires 213k parameters, while their compact model (6.4k params) reaches 88.0\%. Our model reaches 90.1\% with only 3.5k parameters, improving compact-model accuracy by 2.1 points while using fewer parameters.

%\begin{table}[htbp]
%\caption{Comparison with Baselines and State-of-the-Art}
%\label{tab:comparison}
%\centering
%\resizebox{\columnwidth}{!}{%
%\begin{tabular}{lccccc}
%\toprule
%\textbf{Approach} & \textbf{Params} & \textbf{Acc.} & \textbf{Size} & \textbf{On-Chip$^{a}$} & \textbf{Mobile} \\
%\midrule
%Traditional CNN$^{*}$ & 250k+ & 89.2\% & 1.2 MB & No & No \\
%Traditional LSTM$^{*}$ & 180k+ & 87.5\% & 0.8 MB & Marginal & No \\
%Manual CNN-LSTM$^{*}$ & 95k+ & 91.3\% & 0.4 MB & Possible & No \\
%Fontaine IQ-CNN \cite{fontaine2019low} & 213k & 97.8\% & 852 KB & No & No \\
%Fontaine Manual RF \cite{fontaine2019low} & 6.4k & 88.0\% & 26 KB & Yes & No \\
%\textbf{NAS (Ours, updated)} & \textbf{3.5k} & \textbf{90.1\%} & %\textbf{14 KB} & \textbf{Yes} & \textbf{In progress} \\
%\bottomrule
%\end{tabular}%
%}

%\begin{flushleft}
%\footnotesize
%$^{*}$ Same dataset, preprocessing, and training. $^{a}$ On-chip SRAM feasibility (typical FPGA: 100--500~KB BRAM).
%\end{flushleft}
%\end{table}

% =========================================================
\section{Discussion and Conclusion}
% =========================================================

Our results demonstrate that evolutionary NAS can discover architectures that strongly favor efficiency while preserving useful classification accuracy. The updated 3,539-parameter model highlights an aggressive compactness regime where the model remains practical for on-chip SRAM deployment and low-compute inference.

%The evaluation uses 135 files across three classes; larger datasets would strengthen generalization claims. It operates on preprocessed segments rather than continuous signals, and hardware resource estimates are analytical projections pending FPGA synthesis.

Our updated evolutionary NAS run discovered an ultra-compact model achieving 90.1\% accuracy with significant parameter reduction (3,539 parameters), enabling spectrum monitoring and cognitive radio under strict edge constraints.

\bibliographystyle{IEEEtran}
\bibliography{References/references}

\end{document}